\documentclass[../../main.tex]{subfiles}
\begin{document}
% SOURCE p=226; printed=211; object=ch07-positive-functionals-title; region=lower
\section{$C_c(X)$ 上的正线性泛函}\label{section:ch07-positive-functionals}

% SOURCE p=226; printed=211; object=ch07-positive-functional-definition; region=bottom
\begin{definition}[正线性泛函]\label{definition:ch07-positive-functional}
回顾 $C_c(X)$ 是 $X$ 上具有紧支集的连续函数所成的空间。$C_c(X)$ 上的线性泛函 $I$ 称为\textbf{正的}，是指只要 $f\geqslant0$，就有 $I(f)\geqslant0$。
\end{definition}
这个定义没有提及连续性，但正性本身便蕴含一种相当强的连续性性质。
% SUPPLEMENT: id=v14-ch07-positive-functional-motivation; parent=ch07-positive-functional-definition; kind=proof-detail+motivation+hypotheses; origin=independent-derivation
\begin{remark}
\textbf{顺序、实值性与动机（重构新增）。}这里可取复连续函数空间和复线性泛函；$f\geqslant0$ 专指 $f$ 是逐点非负的实值函数，而 $I(f)\geqslant0$ 要求其值为非负实数。若 $h$ 为实值紧支连续函数，则 $h^+,h^-\in C_c(X)$，所以 $I(h)=I(h^+)-I(h^-)$ 为实数。若实值 $h\leqslant k$，正性又给出 $I(h)\leqslant I(k)$。这正是积分的单调性所具有的性质；后面的表示定理将从这项顺序性质恢复测度。定义域是向量空间，使这些差及线性组合合法；紧支集为后续的局部范数界提供一个固定紧集。此处的函数取实际点值，尚未对某个测度作几乎处处商化。
\end{remark}
% END SUPPLEMENT: id=v14-ch07-positive-functional-motivation

% SOURCE p=227; printed=212; object=ch07-positive-functional-local-bound; number=7.1; region=top
\begin{proposition}\label{proposition:ch07-positive-local-bound}
若 $I$ 是 $C_c(X)$ 上的正线性泛函，则对每个紧集 $K\subseteq X$，存在常数 $C_K$，使所有满足 $\operatorname{supp}(f)\subseteq K$ 的 $f\in C_c(X)$ 都有
\begin{align*}
|I(f)|\leqslant C_K\|f\|_u.
\end{align*}
\end{proposition}
% SOURCE p=227; printed=212; object=ch07-positive-functional-local-bound-proof; region=upper
\begin{proof}
只须考虑实值 $f$。给定紧集 $K$，由 Urysohn 引理选取 $\phi\in C_c(X,[0,1])$，使 $\phi=1$ 于 $K$。若 $\operatorname{supp}(f)\subseteq K$，则 $|f|\leqslant\|f\|_u\phi$，即 $\|f\|_u\phi-f\geqslant0$ 且 $\|f\|_u\phi+f\geqslant0$。因此
\begin{align*}
\|f\|_u I(\phi)-I(f)\geqslant0,
\qquad \|f\|_u I(\phi)+I(f)\geqslant0,
\end{align*}
从而 $|I(f)|\leqslant I(\phi)\|f\|_u$。
\end{proof}
% SUPPLEMENT: id=v14-ch07-positive-local-bound-details; parent=ch07-positive-functional-local-bound; kind=proof-detail+motivation+hypotheses; origin=independent-derivation
\begin{remark}
\textbf{复值情形与逐条件说明（重构新增）。}所用 Urysohn 引理是\cref{theorem:ch04-local-urysohn}，取其中的开集为 $X$。若 $I(f)\ne0$，选 $|c|=1$ 使 $cI(f)=|I(f)|$；若 $I(f)=0$ 则结论已成立。由实函数的实值性及复线性，
\begin{align*}
|I(f)|=\operatorname{Re}I(cf)=I(\operatorname{Re}(cf))
\leqslant I(\phi)\|f\|_u,
\end{align*}
因为 $\operatorname{Re}(cf)\leqslant|f|\leqslant\|f\|_u\phi$。所以同一个 $C_K=I(\phi)$ 也控制全部复值函数。若 $K=\varnothing$，只有零函数，取 $C_K=0$ 即可。

动机是用一个固定的非负截断函数同时支配紧集上的所有函数；正性把逐点比较转换为泛函比较。紧性与 LCH 条件保证能选这样的紧支截断；固定 $K$ 使常数不随 $f$ 改变。结论特别说明：支集都在同一紧集的函数列若一致收敛到零，其泛函值也趋零。但它不声称在整个 $C_c(X)$ 上有一个统一常数。譬如实线上 Lebesgue 积分是正泛函，取越来越宽的平台函数，其一致范数为一而积分可任意大。
\end{remark}
% END SUPPLEMENT: id=v14-ch07-positive-local-bound-details

% SOURCE p=227; printed=212; object=ch07-integral-positive-functional; region=upper-middle
若 $\mu$ 是 $X$ 上的 Borel 测度，并且对每个紧集 $K\subseteq X$ 有 $\mu(K)<\infty$，则显然 $C_c(X)\subseteq L^1(\mu)$，故映射 $f\mapsto\int f\,\mathrm d\mu$ 是 $C_c(X)$ 上的正线性泛函。本节的主要结果是，每个正线性泛函都以这种方式产生；此外，可以对 $\mu$ 加上某些正则性条件，使其唯一。这些条件如下。
% SUPPLEMENT: id=v14-ch07-integral-functional-details; parent=ch07-integral-positive-functional; kind=proof-detail+motivation+hypotheses; origin=independent-derivation
\begin{remark}
\textbf{可积性与正性的核验（重构新增）。}连续函数 Borel 可测；若 $f\in C_c(X)$，其紧支集记为 $K$，则 $f$ 有界，并有 $\int|f|\,\mathrm d\mu\leqslant\|f\|_u\mu(K)<\infty$。因而积分的线性性及非负函数积分非负给出所述泛函。这里 $C_c(X)\subseteq L^1$ 按通常约定指把每个实际函数送到其几乎处处等价类后可积；这张映射未必单射。紧集有限足以控制每一个输入，不要求 $\mu(X)$ 有限。要从积分值反推所有 Borel 集的测度，则还需要下面的拓扑逼近条件。
\end{remark}
% END SUPPLEMENT: id=v14-ch07-integral-functional-details

% SOURCE p=227; printed=212; object=ch07-regular-measure-definition; region=middle
\begin{definition}[外正则、内正则与正则测度]\label{definition:ch07-regular-measure}
设 $\mu$ 为 $X$ 上的 Borel 测度，$E\in\mathscr B_X$。若
\begin{align*}
\mu(E)=\inf\{\mu(U):E\subseteq U,\ U\text{ 开}\},
\end{align*}
则称 $\mu$ 在 $E$ 上\textbf{外正则}；若
\begin{align*}
\mu(E)=\sup\{\mu(K):K\subseteq E,\ K\text{ 紧}\},
\end{align*}
则称 $\mu$ 在 $E$ 上\textbf{内正则}。若 $\mu$ 在所有 Borel 集上都外正则且内正则，则称 $\mu$ \textbf{正则}。
\end{definition}
当 $X$ 不是 $\sigma$ 紧空间时，要求正则性可能过强，因此采用以下定义。
% SOURCE p=227; printed=212; object=ch07-radon-definition; region=middle-lower
\begin{definition}[Radon 测度]\label{definition:ch07-radon-measure}
$X$ 上的 \textbf{Radon 测度}是满足下列三项条件的 Borel 测度：在所有紧集上有限，在所有 Borel 集上外正则，在所有开集上内正则。
\end{definition}
我们将在\cref{section:ch07-regularity-approximation}证明，Radon 测度在其所有 $\sigma$ 有限集上也内正则。
% SUPPLEMENT: id=v14-ch07-regular-radon-motivation; parent=ch07-radon-definition; kind=proof-detail+motivation+hypotheses; origin=independent-derivation
\begin{remark}
\textbf{逼近方向与逐条件说明（重构新增）。}单调性总给出 $\mu(K)\leqslant\mu(E)\leqslant\mu(U)$，所以这两个等式分别要求外面的开集和里面的紧集能达到任意精度。若 $\mu(E)<\infty$，外正则等价于每个 $\varepsilon>0$ 有开集 $U\supseteq E$ 满足 $\mu(U\setminus E)<\varepsilon$；内正则等价于可取紧集 $K\subseteq E$ 使 $\mu(E\setminus K)<\varepsilon$。这些等价性来自可测不交差的可加性。若 $\mu(E)=\infty$，外正则由单调性自动成立，而内正则要求紧子集的测度无上界；不能在此用 $\infty-\infty$ 表示误差。

紧集在 Hausdorff 空间中是 Borel 集，故上述内逼近合法。紧集有限使紧支函数可积；开集内正则使开集测度可从紧支连续函数恢复；全部 Borel 集外正则再把开集上的信息扩展到全部 Borel 集。三项在表示定理中承担不同任务。空集可用自身作为开集和紧集，两个取界集合始终非空（外侧还总可取 $X$）。本书“Radon”采用这里的三项约定；不要在未证明的情况下额外加入全体 Borel 集内正则。LCH 条件下，紧集有限也给出每点的某个开邻域测度有限：取一个紧邻域，其内部即可。
\end{remark}
% END SUPPLEMENT: id=v14-ch07-regular-radon-motivation

% SOURCE p=227; printed=212; object=ch07-subordinate-function-notation; region=lower
\begin{definition}[记号 $f\prec U$]\label{definition:ch07-subordinate-function}
若 $U$ 在 $X$ 中开且 $f\in C_c(X)$，写 $f\prec U$ 表示 $0\leqslant f\leqslant1$ 且 $\operatorname{supp}(f)\subseteq U$。这比 $0\leqslant f\leqslant\chi_U$ 略强；后者只蕴含 $\operatorname{supp}(f)\subseteq\overline U$。
\end{definition}
% SUPPLEMENT: id=v14-ch07-subordinate-function-motivation; parent=ch07-subordinate-function-notation; kind=proof-detail+example+motivation+hypotheses; origin=independent-derivation
\begin{remark}
\textbf{支集条件的意义（重构新增）。}支集是非零点集的闭包，因此即使函数在 $U^c$ 上处处为零，非零点的极限仍可能落在 $\partial U$。例如在 $\mathbb R$ 上取 $U=(0,1)$、$f(x)=\max(0,1-|2x-1|)$，则 $0\leqslant f\leqslant\chi_U$，但支集为 $[0,1]$，故 $f\not\prec U$。要求整个紧支集位于 $U$ 内，才能对该支集选取 $U$ 的有限开子覆盖并使用单位分解；界 $f\leqslant1$ 使 $I(f)$ 扮演“集合大小”的下界。$f=0$ 总满足 $f\prec U$，因此后面取上确界的集合不会为空。
\end{remark}
% END SUPPLEMENT: id=v14-ch07-subordinate-function-motivation
\subfile{section-01-riesz.tex}
\subfile{section-01-exercises.tex}
\end{document}
